\section{Discussion}
\label{sec:discussion}

\subsection{Hypothesis Assessment}

The hypothesis is \textbf{confirmed with qualification}. A Go SMT
with Poseidon2 achieves state root computation under 50~ms for up
to 250 transactions per block on a 10,000-entry tree at depth 32.
However, the qualification is significant: at production EVM depths
(160--256), the 50~ms target becomes considerably more challenging.

\textbf{Prediction reconciliation:}

\begin{table}[h]
\centering
\caption{Prediction vs. observed results}
\label{tab:predictions}
\begin{tabular}{llll}
\toprule
\textbf{Prediction} & \textbf{Expected} & \textbf{Observed} & \textbf{Status} \\
\midrule
Hash speedup vs.\ JS & 50--200$\times$ & 12.6$\times$ & Below range \\
Insert latency & $< 100$~$\mu$s & 125.66~$\mu$s & Marginal miss \\
Proof generation & $< 10$~$\mu$s & 10.06~$\mu$s & At boundary \\
Batch 10K $< 50$~ms & $< 50$~ms & 46.05~ms (250 tx) & PASS \\
Memory $< 500$~MB & $< 500$~MB & 29.3~MB & PASS \\
\bottomrule
\end{tabular}
\end{table}

The hash speedup prediction (50--200$\times$) was overly optimistic.
The 12.6$\times$ observed speedup reflects the narrower gap between
V8 BigInt and Go assembly than anticipated, as well as the different
round counts between Poseidon and Poseidon2. Insert latency at
125.66~$\mu$s narrowly missed the 100~$\mu$s target but remains
14.3$\times$ faster than the TypeScript baseline.

\subsection{Depth Sensitivity Analysis}

The cost model (Equation~\ref{eq:cost}) predicts that operation
latency scales linearly with tree depth. At depth 32, batch cost is
dominated by $32 \times C_{\text{hash}}$ per update. At depth 160
(Ethereum address length) or 256 (full key space), operations will
be 5--8$\times$ slower:

\begin{table}[h]
\centering
\caption{Projected batch latency at production depths}
\label{tab:depth-projection}
\begin{tabular}{lrrr}
\toprule
\textbf{Batch size} & \textbf{Depth 32} & \textbf{Depth 160 (est.)} & \textbf{Depth 256 (est.)} \\
\midrule
10 tx  & 1.86~ms  & 9.3~ms   & 14.9~ms \\
50 tx  & 9.03~ms  & 45.2~ms  & 72.2~ms \\
100 tx & 18.77~ms & 93.9~ms  & 150.2~ms \\
250 tx & 46.05~ms & 230.3~ms & 368.4~ms \\
\bottomrule
\end{tabular}
\end{table}

At depth 160, the 50~ms target is achievable only for blocks of
$\leq 50$ transactions. At depth 256, even 50-transaction blocks
exceed the budget (72.2~ms projected). This analysis identifies
three mitigation strategies:

\begin{enumerate}
  \item \textbf{Compact SMT:} Use extension nodes (as in Jellyfish
    Merkle Tree~\cite{aptos2024jellyfish}) to skip sparse regions,
    reducing effective depth to $O(\log n)$ where $n$ is the number
    of populated leaves.
  \item \textbf{Batch parallelism:} Updates to disjoint subtrees
    can be parallelized across goroutines, achieving near-linear
    speedup for non-overlapping key ranges.
  \item \textbf{Incremental root computation:} Cache intermediate
    nodes and only recompute the path from modified leaves to root,
    avoiding redundant hash computations for unchanged subtrees.
\end{enumerate}

All three techniques are employed by production systems: Polygon
zkEVM uses sparse node pruning~\cite{polygon2024zkprover}, zkSync
Era uses incremental updates~\cite{zksync2024boojum}, and Scroll
uses a modified trie with path compression~\cite{scroll2024architecture}.

\subsection{Library Selection: gnark-crypto vs. vocdoni/arbo}

The experiment used gnark-crypto's Poseidon2 for maximum native
performance. However, the library choice has significant downstream
implications:

\begin{table}[h]
\centering
\caption{Go SMT library comparison}
\label{tab:libraries}
\begin{tabular}{lll}
\toprule
\textbf{Criterion} & \textbf{gnark-crypto} & \textbf{vocdoni/arbo} \\
\midrule
Hash function   & Poseidon2        & Poseidon (original) \\
Performance     & Higher (asm)     & Lower ($\sim$2$\times$) \\
Circuit compat. & gnark/halo2      & circom/snarkjs \\
Persistent DB   & None (custom)    & BadgerDB/Pebble \\
Compact SMT     & No               & Yes (built-in) \\
Maturity        & Production       & Production \\
\bottomrule
\end{tabular}
\end{table}

\textbf{Critical observation:} Poseidon2 produces \textit{different
hashes} than circomlibjs Poseidon. The two are not cross-compatible.
The choice of hash function in the state database must align with the
prover's circuit library:
\begin{itemize}
  \item If the prover uses gnark or halo2 (Rust/Go native proving),
    gnark-crypto with Poseidon2 is optimal.
  \item If the prover uses circom/snarkjs (as in the current Validium
    MVP), vocdoni/arbo with original Poseidon is required for
    hash compatibility.
\end{itemize}

This decision propagates to every downstream component and must be
resolved at the architecture level before implementation begins.

\subsection{MPT vs. SMT Trade-offs}

While our results validate SMT for the Basis Network use case, the
MPT vs.\ SMT decision involves trade-offs beyond raw performance:

\textbf{SMT advantages:}
\begin{itemize}
  \item Fixed-depth proofs simplify ZK circuit design.
  \item Uniform proof structure for inclusion and exclusion.
  \item ZK-friendly hash functions (Poseidon/Poseidon2) reduce
    in-circuit constraint count by 500$\times$ vs.\ Keccak.
  \item Deterministic proof size enables predictable proving costs.
\end{itemize}

\textbf{MPT advantages:}
\begin{itemize}
  \item Native Ethereum compatibility (no state translation layer).
  \item Dynamic depth adapts to actual key distribution.
  \item Mature tooling in go-ethereum ecosystem.
  \item Direct compatibility with Ethereum JSON-RPC state proofs.
\end{itemize}

For the Basis Network architecture---where each enterprise operates
a dedicated chain and EVM compatibility is provided at the execution
layer, not the state layer---SMT with Poseidon2 provides the optimal
balance of ZK efficiency and performance.

\subsection{Reconciliation with Production Systems}

Our measurements are consistent with published production benchmarks:

\begin{itemize}
  \item \textbf{Polygon zkEVM:} Reports 2--3~ms per Merkle proof
    verification in the zkProver~\cite{polygon2024zkprover}, compared
    to our 0.151~ms for local verification. The difference reflects
    Polygon's depth-256 tree and C++ prover overhead.
  \item \textbf{Scroll:} Uses Go-based state management with Poseidon
    hashing~\cite{scroll2024architecture}. Their zktrie achieves
    comparable per-operation latency to our measurements.
  \item \textbf{zkSync Era:} Uses Rust with Blake2s
    hashing~\cite{zksync2024boojum}, optimizing for native speed
    over ZK-friendliness (re-proving the hash in-circuit).
\end{itemize}

The 10--14$\times$ speedup over TypeScript is consistent with the
Go-vs-JavaScript performance gap observed in production systems
where gnark-based provers outperform snarkjs by similar
margins~\cite{chaliasos2024analyzing}.

\subsection{Implications for Downstream Pipeline}

The results establish the following constraints for the Architect
(implementation) and Prover (verification) stages:

\begin{enumerate}
  \item The state database implementation must use compact SMT with
    path compression for production depths ($\geq 160$).
  \item Hash function choice (Poseidon2 vs.\ original Poseidon)
    must be decided jointly with the prover framework selection.
  \item Persistent storage (LevelDB or Pebble) is required for
    deployments exceeding 100,000 accounts.
  \item Batch parallelism should be investigated for blocks
    exceeding 100 transactions.
  \item Memory budget of 2.9~KB per account should inform
    hardware sizing recommendations.
\end{enumerate}
